\section{Projet 13 — Penalty !}
\setcounter{UPSTInumeroCahier des charges}{0}

\subsection*{Objectif du jeu}
Simuler une séance de tirs au but entre un joueur physique (joueur 1) et un joueur virtuel (joueur 2, le programme). Droite et gauche sont vues du côté du tireur.

\subsection*{Règles du jeu}
\begin{itemize}
  \item Le joueur 1 choisit un nom pour son équipe et un nom pour l'équipe du joueur 2.
  \item Le tireur choisit son tir sous la forme d'un caractère :
  \begin{itemize}
    \item \texttt{'d'} / \texttt{'g'} : à droite / à gauche du gardien, frappe molle et mal placée ;
    \item \texttt{'D'} / \texttt{'G'} : à droite / à gauche du gardien, frappe forte placée petit filet ;
    \item \texttt{'m'} : au milieu, frappe molle à ras de terre ; \texttt{'M'} : au milieu, frappe type panenka.
  \end{itemize}
  \item Les tirs \texttt{'D'}, \texttt{'G'} ou \texttt{'M'} ont $n$ chances sur 8 de terminer \textbf{en dehors de la cage}, où $n$ est le nombre de fois que \texttt{'D'}, \texttt{'G'} ou \texttt{'M'} ont déjà été choisis par ce joueur lors de la séance (tirage aléatoire par le programme).
  \item Le gardien plonge : \texttt{'d'} (à droite), \texttt{'g'} (à gauche) ou \texttt{'m'} (une partie du corps reste au milieu).
  \item Le gardien gagne si le tir est hors cadre, ou si le tir est \texttt{'d'}, \texttt{'g'} ou \texttt{'m'} et que le gardien a choisi respectivement \texttt{'d'}, \texttt{'g'} ou \texttt{'m'}. Sinon le tireur marque.
  \item Les choix du joueur virtuel sont aléatoires. Les joueurs changent de rôle après chaque tir.
  \item Séance de \textbf{5 penalties} tirés en alternance ; en cas d'égalité, \textbf{mort subite}. Le jeu s'arrête dès que l'issue est acquise (ex.\ 4 à 2 après 4 tirs).
\end{itemize}

\subsection*{Cahier des charges}
\begin{UPSTICahierDesCharges}{1}
\begin{itemize}
  \item Séance complète joueur physique contre programme : saisie des noms, choix des tirs et des plongeons, gestion de la probabilité de tir hors cadre, score.
  \item Fin anticipée de la séance, mort subite, annonce du vainqueur.
\end{itemize}
\end{UPSTICahierDesCharges}

\begin{UPSTICahierDesCharges}{2}
\begin{itemize}
  \item Ranger le résultat de chaque tir dans un tableau (un par joueur, 20 cases pour prévoir la mort subite) : 0 = manqué, 1 = arrêté par le gardien, 2 = réussi.
  \item En fin de partie, afficher le résumé de la séance \textbf{en couleur}, par exemple :
\begin{verbatim}
Penalty 1 pour OM : Arrêté par le gardien
Penalty 2 pour OM : Réussi
...
Penalty 5 pour PSG : Réussi

PSG win !!!
\end{verbatim}
\end{itemize}
\end{UPSTICahierDesCharges}

\begin{UPSTICahierDesCharges}{3}
\begin{itemize}
  \item Les 2 joueurs sont des joueurs physiques (choix du mode par menu).
\end{itemize}
\end{UPSTICahierDesCharges}
